\chapter{Unpacking and pulse-shape analysis}
\label{ch:unpacking}

This chapter covers the first two stages of the reconstruction of experimental data: turning the
files written by the data-acquisition systems into \code{AtRawEvent} objects (\emph{unpacking}), and
turning the digitized pad traces of an \code{AtRawEvent} into three-dimensional \code{AtHit} objects
collected in an \code{AtEvent} (\emph{pulse-shape analysis}, PSA). Between the two sit optional
trace filters (baseline, fixed-pattern noise, Fourier filters) and event-selection tasks. The same
PSA stage also runs on simulated data, after the digitization produces an \code{AtRawEvent} from
Monte Carlo points; Section~\ref{sec:unp-pitfalls} explains why the PSA configuration must then be
the same for data and simulation.

The code lives in three places of the repository:
\begin{itemize}
  \item \url{AtUnpack/}: the unpacker interface, the GET (GRAW and HDF5) unpackers, the FRIBDAQ
        auxiliary-stream unpackers, the FRIBDAQ scaler reader and the FPN pedestal helper.
  \item \url{AtReconstruction/AtFilter/} and the tasks \code{AtFilterTask}, \code{AtAuxFilterTask},
        \code{AtDataReductionTask}, \code{AtCopyTreeTask}, \code{AtCopyAuxTreeTask},
        \code{AtDataCleaningTask} in \url{AtReconstruction/}.
  \item \url{AtReconstruction/AtPulseAnalyzer/}: \code{AtPSA} and its implementations, run by
        \code{AtPSAtask} (\url{AtReconstruction/AtPSAtask.h}).
\end{itemize}

A typical data chain, with the default branch names, is
\begin{verbatim}
  .h5/.graw    --AtUnpackTask-->  AtRawEvent
               --AtFilterTask (optional)-->  AtRawEventFiltered
               --AtPSAtask-->  AtEventH
               --space-charge correction, AtDataCleaningTask (optional)--> ...
\end{verbatim}

%=====================================================================================================
\section{The unpacker framework}
\label{sec:unp-framework}

\subsection{\code{AtUnpacker}}

\code{AtUnpacker} (\url{AtUnpack/AtUnpacker.h}) is the abstract interface every unpacker implements.
It holds a shared pointer to the pad map (\code{AtMap}), the input file name and two event
counters:
\begin{itemize}
  \item \code{fEventID}: the internal event number, starting at the value given to
        \code{SetInitialEventID()} (default 0). This is the number written into the
        \code{AtRawEvent} with \code{SetEventID()}.
  \item \code{fDataEventID}: the event number as the DAQ (or merger) recorded it. Each unpacker
        computes it from \code{fEventID} and the first event found in the file.
\end{itemize}
The pure virtual methods are \code{Init()}, \code{FillRawEvent(AtRawEvent\&)},
\code{IsLastEvent()} and \code{GetNumEvents()}. \code{SetSaveFPN(bool)} (default \code{false})
asks the unpacker to store the fixed-pattern-noise (FPN) channels of the AGET chips in the event
as well (\code{AtRawEvent::AddFPN}).

\subsection{\code{AtUnpackTask}}

\code{AtUnpackTask} (\url{AtUnpack/AtUnpackTask.h}) is the \code{FairTask} that owns an unpacker,
passed to the constructor as a \code{std::unique\_ptr<AtUnpacker>}. In \code{Init()} it calls the
unpacker's \code{Init()} and registers the output branch; in \code{Exec()} it fills one
\code{AtRawEvent} per call.

\begin{center}
\begin{tabular}{p{5.2cm}p{3.0cm}p{6.5cm}}
\toprule
Setter & Default & Meaning \\
\midrule
\code{SetOuputBranchName(name)} & \code{"AtRawEvent"} & Output branch (note the spelling of the
setter). \\
\code{SetPersistence(bool)} & \code{true} & Write the raw event to the output file. \\
\code{GetNumEvents()} & -- & Forwards to the unpacker; use it as the upper limit of
\code{FairRunAna::Run(0, n)}. \\
\bottomrule
\end{tabular}
\end{center}

The unpacking runs without a FairRoot source, so the run length must be given explicitly. All the
macros in \url{macro/Unpack_HDF5/} use the same pattern:
\begin{verbatim}
  run->Init();
  auto numEvents = unpackTask->GetNumEvents();
  run->Run(0, numEvents);
\end{verbatim}
After the unpacker reports its last event, further \code{Exec()} calls log
\code{"Hit last event at: ..."} and mark the event bad. \code{AtUnpackTask::FinishEvent()}, which
would stop the run, carries the comment that it is normally not called because there is no source.

%=====================================================================================================
\section{GET data in HDF5 files}
\label{sec:unp-hdf5}

\subsection{\code{AtHDFUnpacker}: the remerged layout}

\code{AtHDFUnpacker} (\url{AtUnpack/AtHDFUnpacker.h}) reads the \emph{remerged} HDF5 layout used by
most AT-TPC experiments at FRIB and Argonne: a group \code{/get} with one dataset
\code{evt<N>\_data} per event (shape \{number of pads, 517\}, native 16-bit integers), an optional
\code{evt<N>\_header} dataset, and a \code{/meta} group.

\paragraph{Event range.} \code{Init()} opens the file and determines the first and last DAQ event
number. If \code{/meta/meta} exists, the first event is word 0 and the last event is word 2 of that
dataset. If the \code{/meta} group lacks the \code{meta} dataset (the code notes that some merger
versions write it that way, for example 12 of 104 a2091 runs), the unpacker logs a warning and
scans the names of all objects in \code{/get} instead, taking the minimum and maximum number
found. \code{GetNumEvents()} returns \code{last - first + 1}; if this differs from the number of
objects in \code{/get}, an error is logged (the count includes the header datasets, so this error
does not by itself mean data are missing). The DAQ number of the event being unpacked is
\code{fDataEventID = first + fEventID}.

\paragraph{One pad.} Each row of \code{evt<N>\_data} has 517 words: the first four are the CoBo,
AsAd, AGET and channel numbers (the fifth identification word is skipped), followed by 512 ADC
samples. The \code{AtPadReference} built from the first four words is resolved through the map:
\begin{enumerate}
  \item if \code{SetSaveFPN(true)} was called and the channel is an FPN channel, it is stored with
        \code{AtRawEvent::AddFPN};
  \item if the map declares the channel auxiliary (\code{AtMap::AddAuxPad}), it is stored as an
        auxiliary pad under the name given in the map;
  \item otherwise it is stored as a pad with the pad number from \code{AtMap::GetPadNum}, and the
        pad centre and size identifier are taken from the map.
\end{enumerate}
The HDF5 unpackers do not consult the inhibit map: pads declared inhibited in \code{AtMap} are
unpacked like any other pad. \code{AtGRAWUnpacker} and the digitization (\code{AtPulse}) do skip
them.

\paragraph{Baseline.} With \code{SetBaseLineSubtraction(true)} the baseline is the mean of the first
20 samples (time buckets 0--19), and \code{fAdc[tb] = raw[tb] - baseline} for all 512 buckets; the
pad is then flagged \code{IsPedestalSubtracted()}. With the default (\code{false}) the baseline is
zero, \code{fAdc} equals the raw ADC, and the pad is \emph{not} flagged; every PSA class then logs
\code{"Pedestal should be subtracted to use this class!"} for each pad. The raw samples are always
kept in \code{fRawAdc}.

\paragraph{Timestamps.} \code{SetNumberTimestamps(n)} (default 0) copies words $1\ldots n$ of
\code{evt<N>\_header} into the event (\code{AtRawEvent::SetTimestamp(value, i)}). The meaning of
each header word is not documented in the code; the a1975 macros use \code{n = 2} for the GET
stream and \code{n = 1} for the FRIBDAQ stream. If the header cannot be read, no timestamps are
set.

\paragraph{Missing events.} If \code{evt<N>\_data} does not exist, the pad count is zero and the
event is still marked good. The a1975 macros therefore define a selection
\code{(evt->GetNumPads() > 0) \&\& evt->IsGood()}.

\subsection{\code{AtMergerHDFUnpacker}: the raw merger layout}

\code{AtMergerHDFUnpacker} reads files written directly by \code{libattpc\_merger} before the
remerge step: a group \code{/events} with one subgroup \code{event\_<N>} per event, the pad data in
\code{event\_<N>/get\_traces} (same 517-word rows, so pad decoding is inherited unchanged), FRIBDAQ
data in \code{event\_<N>/frib\_physics/1903} and \code{.../977}, and scalers in \code{/scalers}. The
event range is the minimum and maximum \code{<N>} among the subgroup names. This layout carries no
GET header, so no timestamps are set. Events missing from the file give an empty event.

The static function \code{AtMergerHDFUnpacker::IsMergerFile(file)} returns \code{true} when the
file has an \code{/events} group and no \code{/get} group, so a macro can choose the reader:
\begin{verbatim}
  std::unique_ptr<AtHDFUnpacker> unp;
  if (AtMergerHDFUnpacker::IsMergerFile(inputFile.Data()))
     unp = std::make_unique<AtMergerHDFUnpacker>(map);
  else
     unp = std::make_unique<AtHDFUnpacker>(map);
\end{verbatim}

\subsection{\code{AtFRIBLinkedHDFUnpacker} and \code{AtFRIBSiUnpacker}}

\code{AtFRIBLinkedHDFUnpacker} (\url{AtUnpack/AtFRIBLinkedUnpacker.h}; note that the file name and
the class name differ) reads files in which GET and FRIBDAQ data were already linked by timestamp,
both under \code{/events/event\_<N>}. It assumes events are numbered contiguously from 0
(\code{last = number of objects - 1}). Pads come from \code{event\_<N>/get\_traces}. The
timestamps are the HDF5 attributes \code{timestamp} and \code{timestamp\_other} of the
\code{get\_traces} object (one or two timestamps are set). For every path in \code{fFribPaths}
(default \code{\{"1903"\}}; \code{SetFribPaths}, \code{AddFribPath}) the dataset
\code{event\_<N>/frib\_physics/<path>} is read and each of its channels is added as an
\code{AtGenericTrace}, with the trace length taken from the dataset shape. These generic traces are
not baseline-subtracted.

\code{AtFRIBSiUnpacker} extends it for files with silicon detectors read by GET: pads are read from
\code{event\_<N>/get/pads}, \code{fFribPaths} defaults to \code{977, 1903, 1904, 1905, 1906}, and
the channels of the datasets \code{si\_downstream\_back}, \code{si\_downstream\_front},
\code{si\_upstream\_back}, \code{si\_upstream\_front} under \code{event\_<N>/get/} are added as
auxiliary pads named \code{Si\_<cobo>\_<asad>\_<aget>\_<ch>}. Test macros:
\url{macro/tests/AT-TPC/run_unpack_linked.C} and \url{macro/tests/AT-TPC/run_unpack_si.C}.

%=====================================================================================================
\section{The FRIBDAQ auxiliary stream}
\label{sec:unp-frib}

At FRIB the ion chamber (IC) and other auxiliary signals are digitized by the FRIBDAQ system (a
SIS3300-type digitizer, module 1903) and stored next to the GET data. Two unpackers read them into
a separate \code{AtRawEvent} whose content is a list of \code{AtGenericTrace} objects
(\code{AtRawEvent::GetGenTraces()}), not pads.

\begin{center}
\begin{tabular}{p{4.4cm}p{5.4cm}p{4.8cm}}
\toprule
Unpacker & Dataset per event & Notes \\
\midrule
\code{AtFRIBHDFUnpacker} & \code{/frib/evt/evt<N>\_1903}, shape \{2048, 8\} & Remerged layout.
Event range from the dataset names, header datasets skipped. \\
\code{AtMergerFRIB-} \code{HDFUnpacker} & \code{/events/event\_<N>/} \code{frib\_physics/1903} & Raw merger layout,
same shape and decoding. \\
\bottomrule
\end{tabular}
\end{center}

Each of the 8 columns becomes one \code{AtGenericTrace} of 2048 samples, with the column index as
trace ID. The header of \code{AtMergerFRIBHDFUnpacker} states that the ion chamber is channel 0, so
the a1975 analysis macros read the IC from \code{GetGenTraces()[0]}. With
\code{SetBaseLineSubtraction(true)} the baseline is the mean of samples 5--24 of the FRIB trace
(the same index range as for GET rows, where the first five words are identifiers; here there are
no identifier words). \code{AtFRIBHDFUnpacker::Init()} logs an error when the number of events
derived from the dataset names is not one third of the number of objects in \code{/frib/evt}
(three objects per event in that layout); \code{AtMergerFRIBHDFUnpacker} does not make that
check because its layout has one subgroup per event.

The FRIBDAQ stream is unpacked in a separate run that writes its own file
named \code{<run>\_FRIB.root}. The example from
\url{macro/Unpack_HDF5/a1975/unpack_a1975_FRIBDAQ.C}:
\begin{verbatim}
  auto unpacker = std::make_unique<AtFRIBHDFUnpacker>(fAtMapPtr);
  unpacker->SetInputFileName(inputFile.Data());
  unpacker->SetNumberTimestamps(1);
  unpacker->SetBaseLineSubtraction(true);
  auto unpackTask = new AtUnpackTask(std::move(unpacker));
  unpackTask->SetPersistence(true);
  run->AddTask(unpackTask);
\end{verbatim}

\subsection{Auxiliary channels in the GET stream}

Auxiliary signals recorded by GET itself (for example a mesh or an IC fed into an AGET channel)
are declared in the map before unpacking:
\begin{verbatim}
  fAtMapPtr->AddAuxPad({cobo, asad, aget, ch}, "IC");
\end{verbatim}
The HDF5 unpackers then store that channel as an auxiliary pad, retrieved with
\code{AtRawEvent::GetAuxPad(name)}, instead of a pad of the pad plane, and the PSA never sees it. The macro
\url{macro/Unpack_HDF5/a1975/UKF/pipeline/discover_aux_a1975.C} scans a run for channels that could
be added this way.

\subsection{How the streams are aligned}
\label{sec:unp-align}

There is no event-by-event linking task for the GET and FRIBDAQ streams in the remerged layout.
Each unpacker numbers its events from 0 (\code{fEventID}) relative to the first event it finds in
its own group, and writes that number into the \code{AtRawEvent}; the DAQ event number
(\code{fDataEventID}) is not stored for the GET stream. The analysis macros pair the two files by
tree entry. For example \url{macro/Unpack_HDF5/a1975/UKF/pd/ex_pd_a1975.C} loops over
\code{std::min(tu->GetEntries(), tc->GetEntries())} entries and reads entry $i$ of both trees. This
is correct only when
\begin{itemize}
  \item both groups start at the same DAQ event number (compare the \code{"Events: X to Y"} lines
        that both unpackers log at \code{Init()});
  \item neither stream has gaps (a missing GET event still produces an entry, see above);
  \item no task removed entries from one of the two trees (\code{AtDataReductionTask} does, see
        Section~\ref{sec:unp-tasks});
  \item both streams were unpacked to the end. Nothing in the unpackers compares the event count
        of one stream with the other, so a FRIBDAQ stream shorter than the GET stream produces a
        short \code{\_FRIB.root} without any warning, and the \code{std::min} loop then silently
        drops the remaining GET events.
\end{itemize}
In the linked layouts (\code{AtFRIBLinkedHDFUnpacker}) the linking is done upstream and the IC
traces are in the same \code{AtRawEvent} as the pads. For experiments with a separate NSCL/FRIB DAQ
recorded in a HiRAEVT tree, \code{AtLinkDAQTask} links the two by the ratio of the time intervals
between consecutive events; it is compiled only when HiRAEVT is found
(\url{AtReconstruction/CMakeLists.txt}).

%=====================================================================================================
\section{FRIBDAQ scalers: \code{AtFribScalers}}
\label{sec:unp-scalers}

\code{AtFribScalers} (\url{AtUnpack/AtFribScalers.h}) reads the FRIBDAQ scaler counters of a run.
None of the unpackers reads them. The header states that the class was validated bit-identical
against Spyral's reader on a1975 run 0016 for all eleven counters.

\code{ReadRun(fileName)} looks for the scaler group \code{frib/scaler} (remerged layout) and then
\code{scalers} (raw merger layout), and reads the datasets \code{scaler0\_data},
\code{scaler1\_data}, \ldots{} until one is missing; this missing dataset is the normal end of the
scan. Of the 32 words of each dataset, the first eleven are read:

\begin{center}
\begin{tabular}{p{0.9cm}p{3.0cm}p{10.1cm}}
\toprule
Index & Enum & Meaning (from the header) \\
\midrule
0 & \code{kClockFree} & time elapsed while the DAQ was running \\
1 & \code{kClockLive} & time for which the DAQ could accept triggers \\
2 & \code{kTriggerFree} & triggers received \\
3 & \code{kTriggerLive} & triggers that produced events \\
4 & \code{kIcSca} & ion-chamber counts (the beam counter) \\
5 & \code{kMeshSca} & mesh signals \\
6, 7 & \code{kSi1Cfd}, \code{kSi2Cfd} & silicon detectors 1 and 2 \\
8 & \code{kSiPM} & unclear (the header says FRIBDAQ's documentation says so) \\
9 & \code{kIcDs} & downscaled ion chamber \\
10 & \code{kIcCfd} & unclear; equal to \code{kIcSca} in every a1975 run checked \\
\bottomrule
\end{tabular}
\end{center}

\paragraph{The counters are incremental.} Each scaler event holds the counts since the previous
one. A run total is therefore the \emph{sum} over all scaler events of the run, which is what
\code{ReadRun} computes, and not the last value. Several runs are combined with \code{Add()}, which
again sums. Never take the difference between the values at the start and end of a run range.
\code{GetLiveFractionClock()} and \code{GetLiveFractionTrigger()} return live/free for the two
pairs of counters; the header notes that they differ by about 2\% and that the right one depends on
what the yield was normalised to.

\paragraph{A run range is not a run list.} The macro
\url{macro/Unpack_HDF5/a1975/frib_scalers_a1975.C} turns the IC counts into a luminosity. Its
documentation records that summing all runs numbered 16--103 instead of the 47 runs used by the
(d,t) analysis gave 483.7 instead of 330.0. The macro therefore takes a run-list file (one run per
line) and prints a warning when it is given a bare range:
\begin{verbatim}
  root -b -q 'frib_scalers_a1975.C(0,0,"/mnt/f/a1975/h5/",2,300,293,
                                   97.1731,0.613,"d2_runs.txt")'
\end{verbatim}
The core of the macro is
\begin{verbatim}
  AtFribScalers total;
  for (int r : runs) {
     AtFribScalers one;
     if (!one.ReadRun(Form("%srun_%04d.h5", h5dir.Data(), r))) {
        ++nMiss;
        continue;
     }
     total.Add(one);
  }
  const double ic = total.Get(AtFribScalers::kIcSca);
\end{verbatim}
Runs without scaler data are counted and reported; the macro states that the result is then a
lower bound.

%=====================================================================================================
\section{GET data in GRAW files}
\label{sec:unp-graw}

\code{AtGRAWUnpacker} (\url{AtUnpack/AtGRAWUnpacker.h}) decodes the binary GRAW files written by the
GET DAQ, using the \code{GETDecoder2} library in \url{AtUnpack/GETDecoder2/}. It was written for
GADGET~II and also serves AT-TPC data; the \code{deprecated/} directory holds its predecessor
\code{AtCore2}.

The constructor \code{AtGRAWUnpacker(map, numGrawFiles = 4)} creates one decoder and one
\code{AtPedestal} per file stream (one per CoBo or AsAd); with more than one stream the data are
treated as \emph{separated} and the streams are decoded in parallel threads. The input file is a
text list of GRAW files. \code{SetInputFileName(list, "file\%i\_")} assigns each line to decoder
$i$ when it contains the formatted string; with the one-argument version the pattern defaults to
\code{"AsAd\%i"} and a warning is printed.

\paragraph{Event alignment.} At \code{Init()} the unpacker reads the first frame of every stream and
starts at the largest of their event IDs (plus \code{SetInitialEventID}). For each event it asks
every decoder for the frame with that ID and compares the event IDs returned; a mismatch logs
\code{"Event IDs don't match between files!"} and marks the event bad. \code{IsLastEvent()} is
true as soon as any stream has no frame for the next ID. \code{GetNumEvents()} returns $-1$ unless
\code{SetCheckNumEvents()} was called, in which case all files are scanned once (or only the first,
with \code{SetMutantOneRun(true)}). \code{SetUseCoboFrame(true)} (default \code{false}) assembles
each event from a full CoBo frame; the header states this is required for multi-AsAd data such as
the December 2014 alpha runs, together with \code{SetPseudoTopologyFrame(mask, check)}.

\paragraph{FPN subtraction.} With \code{SetSubtractFPN(true)} (the default) each pad is corrected
with the FPN channel of the same AGET quarter (\code{AtMap::GetNearestFPN}: channel 11, 22, 45 or
56). \code{AtPedestal::SubtractPedestal} looks for a window of 10 time buckets, starting at bucket
3 and stepping by 10, whose RMS is below \code{SetFPNSigmaThreshold} (default 5 ADC). It takes the
difference between the means of the pad and of the FPN channel in that window, and sets
\code{adc = raw - (fpn - offset)}. If no window satisfies the RMS cut, it prints a message and
leaves the pad unchanged. Inhibited pads (\code{AtMap::InhibitType} other than \code{kNone}) are
not unpacked.

\paragraph{Baseline subtraction.} If FPN subtraction is off and
\code{SetBaseLineSubtraction(true)}, the baseline is the mean of raw samples 5--24.

\paragraph{Other options.} \code{SetSaveFPN()} stores the FPN channels; \code{SetSaveLastCell(true)}
stores the SCA last-read cell of each AGET as a pad augment \code{"lastCell"}, needed by
\code{AtSCACorrect}. \code{SetIsPositivePolarity()} stores a flag that is not used anywhere in
\code{AtGRAWUnpacker.cxx}: the FPN subtraction always uses the default (positive) polarity of
\code{AtPedestal::SubtractPedestal}.

\code{AtROOTUnpacker} reads SpecMAT-style ROOT files through the same interface; it throws if
\code{SetSaveFPN} is called. The test macro \url{macro/tests/AT-TPC/run_unpack_graw.C} shows a
36-stream GRAW configuration.

%=====================================================================================================
\section{Trace filters}
\label{sec:unp-filters}

\subsection{\code{AtFilter} and \code{AtFilterTask}}

A trace filter implements \code{AtFilter} (\url{AtReconstruction/AtFilter/AtFilter.h}):
\code{Init()} once, \code{InitEvent(rawEvent)} at the start of each event, \code{Filter(pad,
padRef)} for each trace, and \code{IsGoodEvent()} at the end. \code{AtFilterTask}
(\url{AtReconstruction/AtFilterTask.h}) copies the input \code{AtRawEvent} to a new branch and
applies the filter to the copy:

\begin{center}
\begin{tabular}{p{4.0cm}p{4.0cm}p{6.6cm}}
\toprule
Setter & Default & Meaning \\
\midrule
\code{SetInputBranch} & \code{"AtRawEvent"} & \\
\code{SetOutputBranch} & \code{"AtRawEventFiltered"} & \\
\code{SetPersistence} & \code{false} & \\
\code{SetFilterPads} & \code{true} & Filter the pads of the pad plane. \\
\code{SetFilterAux} & \code{false} & Filter the auxiliary pads. \\
\code{SetFilterFPN} & \code{false} & Filter the stored FPN channels. \\
\bottomrule
\end{tabular}
\end{center}

The task takes a raw \code{AtFilter*}. Remember to point the next task (normally \code{AtPSAtask})
at \code{"AtRawEventFiltered"}; otherwise it reads the unfiltered branch. The filtered event is
marked good only if the input event was good and \code{IsGoodEvent()} returns \code{true}.

\subsection{Available filters}

\begin{center}
\begin{tabular}{p{3.9cm}p{10.9cm}}
\toprule
Class & What it does \\
\midrule
\code{AtFilterSubtraction} & Per event, averages the traces of AGET channel 0 of every AsAd whose
samples all stay below \code{SetThreshold} (default 0), and subtracts that average from every pad
of the same CoBo and AsAd. Pads not flagged pedestal-subtracted are left unchanged. AsAds without a
usable channel 0 are logged; with \code{SetIsGood(true)} (default \code{false}) such an event is
marked bad. Constructor \code{(map, numCoBos = 10, numAget = 4)}; the second size is used as the
number of AsAds per CoBo. \\
\code{AtFilterFPN} & Same subtraction, but the reference is the average of the stored FPN channels
(it reads \code{GetFpnPads()}, so the unpacker must be run with \code{SetSaveFPN}). An FPN trace is used only if no sample-to-sample jump exceeds the threshold.
With \code{averageAgets = false} each AGET gets its own reference. \\
\code{AtSCACorrect} & Subtracts a stored average baseline and a phase (last-cell) effect read from
an \code{AtRawEvent} in a file; \code{DoBaseline}, \code{DoPhase} (both set \code{true} by the constructors),
\code{UseChannelZero()}. \\
\code{AtFilterZero} & Fills isolated zero samples by interpolation between neighbours, or copies the
previous sample when the neighbours differ by more than \code{SetThreshold} (default 1000). \\
\code{AtFilterFFT} & Multiplies the Fourier transform of each trace by factors defined on frequency
ranges (\code{AddFreqRange}); \code{SetLowPass(order, cutoff)} builds a Butterworth low-pass.
\code{SetSaveTransform} keeps the transform. \\
\code{AtFilterCalibrate} & Per-pad linear gain calibration (slope and intercept) from
\code{SetCalibrationFile}. \\
\code{AtTrapezoidFilter} & Trapezoidal shaping (\code{SetM}, \code{SetRiseTime},
\code{SetTopTime}, \code{SetDiscriminatorThreshold}). \\
\code{AtRemovePulser} & Removes a pulser signal between two thresholds given to the
constructor. \\
\code{AtFilterDivide}, \code{AtFilterNoOp} & Example filters. \\
\bottomrule
\end{tabular}
\end{center}

\code{AtCalibration} in \url{AtReconstruction/AtPulseAnalyzer/} (gain and jitter files) is used
only by the deprecated \code{AtPSASimple2}; its header recommends \code{AtFilter} instead.

\subsection{\code{AtAuxFilterTask}}

\code{AtAuxFilterTask(filter)} applies a filter to named auxiliary pads (\code{AddAuxPad(name)})
of the input event (\code{SetInputBranchName}, default \code{"AtRawEvent"}). The filtered trace is
added to the \emph{same} event as a new auxiliary pad called \code{<name>Filtered}, and the event's
good flag is replaced by \code{IsGoodEvent()} of the filter.

%=====================================================================================================
\section{Event-selection and tree-copy tasks}
\label{sec:unp-tasks}

\paragraph{\code{AtDataReductionTask}.} It takes a function returning \code{bool}. When the
function returns \code{false}, the task sets \code{IsGood(false)} on the input branch and on every
branch given with \code{SetOutputBranch}, and calls \code{FairRunAna::MarkFill(false)}: the event
is \emph{not written} to the output tree. Overloads accept \code{std::function<bool()>},
\code{std::function<bool(AtBaseEvent*)>} or a templated \code{std::function<bool(T*)>}. The input
branch defaults to \code{"AtRawEvent"}. An example is \url{macro/Unpack_HDF5/e18008_S800/unpack_new.C}:
\begin{verbatim}
  AtDataReductionTask *reduceTask = new AtDataReductionTask();
  reduceTask->SetInputBranch("AtRawEvent");
  reduceTask->SetReductionFunction(&reduceFunc);
\end{verbatim}
Because the rejected entries disappear, a file reduced this way no longer matches a separate
FRIBDAQ file entry by entry.

\paragraph{\code{AtCopyTreeTask} and \code{AtCopyAuxTreeTask}.} \code{AtCopyTreeTask} copies all
branches of the input FairRoot tree to the output; its header states it must be the first task of
the run. It is meant to be combined with \code{AtDataReductionTask} to apply cuts to existing files.
\code{AtCopyAuxTreeTask} copies the matching entries of a tree from another DAQ (\code{SetInputTree},
\code{AddInputFile}, \code{SetOutputFile}), keeping only events whose check branch
(\code{SetCheckEventBranch}, default \code{"AtRawEvent"}) is good, so that the two files stay in
step after a reduction.

%=====================================================================================================
\section{Pulse-shape analysis}
\label{sec:unp-psa}

\subsection{\code{AtPSAtask}}

\code{AtPSAtask} (\url{AtReconstruction/AtPSAtask.h}) reads an \code{AtRawEvent}, copies the event
header into a new \code{AtEvent}, and passes both to the PSA object.

\begin{center}
\begin{tabular}{p{5.0cm}p{3.0cm}p{6.8cm}}
\toprule
Setter & Default & Meaning \\
\midrule
\code{SetInputBranch} & \code{"AtRawEvent"} & Use \code{"AtRawEventFiltered"} after a filter. \\
\code{SetOutputBranch} & \code{"AtEventH"} & \\
\code{SetPersistence} & \code{false} & \\
\code{SetSimlulatedPointBranch} & \code{"AtTpcPoint"} & If this branch exists, MC points are
attached to hits (simulation). \\
\bottomrule
\end{tabular}
\end{center}

The preferred constructor takes a \code{std::unique\_ptr<AtPSA>}. The constructor taking a raw
\code{AtPSA*} is marked deprecated and stores a \emph{clone} of the object: every setter must be
called before the task is constructed, because later calls modify the original, not the copy used
by the task. Events flagged bad are copied to the output without hits.

\subsection{The base class \code{AtPSA}}

\code{AtPSA::Init()} reads from \code{AtDigiPar}: \code{GetTBTime()} (the time-bucket length in ns,
derived from \code{SamplingRate}: 3, 6, 12, 25, 50, 100 MHz give 320, 160, 80, 40, 20, 10 ns, any
other value gives $-1$), \code{DriftVelocity} [cm/\textmu s], \code{ZPadPlane} [mm],
\code{TBEntrance}, \code{BField}, \code{EField} and \code{ReverseDrift}. The number of time
buckets \code{fNumTbs} is fixed at 512. The values are printed at initialisation; check them in
the log.

\paragraph{From time bucket to z.} \code{CalculateZGeo(tb)} returns the drift length measured from
the pad plane:
\begin{equation}
  z = Z_\mathrm{pad} - (\mathrm{TB}_\mathrm{ent} - \mathrm{tb})\,
      \frac{t_\mathrm{TB}\,v_d}{100},
\end{equation}
with $t_\mathrm{TB}$ in ns, $v_d$ in cm/\textmu s and $z$ in mm. A hit at the entrance time
bucket $\mathrm{TB}_\mathrm{ent}$ is at $z = Z_\mathrm{pad}$ (the beam entrance, far from the pad
plane); $z = 0$ is the pad plane. With \url{parameters/ATTPC.a1954.par} ($t_\mathrm{TB} = 160$~ns,
$v_d = 1.30$~cm/\textmu s, $\mathrm{TB}_\mathrm{ent} = 480$, $Z_\mathrm{pad} = 1000$~mm) one time
bucket is 2.08~mm. The fitters convert to the laboratory frame with $z_\mathrm{lab} = Z_\mathrm{pad}
- z$. When \code{ReverseDrift} is set in the parameter file (beam entering through the pad plane),
\code{CalculateZGeo} returns $Z_\mathrm{pad} - z$ so that downstream code needs no change. The
quantity \code{fTB0} printed as ``Pad plane TB'' is not used in the conversion, and the parameter
\code{TB0} of the file is marked deprecated.

\paragraph{\code{CalibrateZ}.} \code{CalibrateZ(tb, padNum)} first adds a per-pad time offset (if
loaded with \code{LoadPadTimeOffsets(csv)}, one value per line in pad order, or
\code{SetPadTimeOffset}) and then calls \code{CalculateZGeo}, unless \code{SetSpyralZ(windowTB,
mmTB = 10, length = 1000)} was called. The latter uses $z = (\mathrm{windowTB} - \mathrm{tb}) /
(\mathrm{windowTB} - \mathrm{mmTB}) \times \mathrm{length}$, which is zero at the \emph{window},
the opposite convention to \code{CalculateZGeo}; the code comment warns that the fitters' $Z_\mathrm{pad}
- z$ then gives a mirrored $z_\mathrm{lab}$. Only \code{AtPSAMax} and \code{AtPSAMultiFit} call
\code{CalibrateZ}; all other PSA classes call \code{CalculateZGeo} directly and ignore per-pad
offsets and \code{SetSpyralZ}.

\paragraph{Thresholds.} \code{SetThreshold(t)} sets the threshold for all pads (and the low
threshold too, unless \code{SetThresholdLow} was called). \code{SetThresholdLow(t)} sets a separate
threshold for pads with size identifier 0 (the small central pads). Both default to $-1$; the
meaning of ``below threshold'' is specific to each class.

\paragraph{Event-level output.} \code{AtPSA::Analyze} loops over the pads, collects the hits,
attaches MC information when present, and fills: the mesh signal (sum of the ADC traces of all pads
that produced at least one hit), the event charge (sum of \code{GetTraceIntegral()} over all
\emph{hits}, so a pad with several hits contributes its trace integral several times when the class
stores the full integral on each hit), the pad multiplicity map and an $r$ variance. Hits are
sorted by \code{AtHit::GetTimeStamp()}, which is an \emph{integer} time bucket.

\paragraph{Per-hit time fields.} \code{AtHit} carries three time fields: \code{fTimeStamp}
(\code{Int\_t}), \code{fTimeStampCorr} and \code{fTimeStampCorrInter} (\code{Double\_t}). Their
content depends on the class, as listed below.

\subsection{\code{AtPSAMax}}

One hit per pad at the maximum of the trace; used by the a1975 data productions.
\begin{itemize}
  \item The maximum is searched in time buckets 20--499. The hit is kept if the maximum exceeds the
        threshold (when the threshold is positive).
  \item Position: pad centre in $x,y$; $z$ = \code{CalibrateZ(maxTB - fPeakingShiftTBs)}, at the
        \emph{integer} time bucket of the maximum. With \code{SetTimeCorrection(true)} (default
        \code{false}) the centre-of-gravity time \code{TBCorr} is used instead.
  \item Charge: the maximum ADC value. Trace integral: the sum of all 512 samples, accumulated in an
        \code{int} (\code{std::accumulate} with an integer initial value), so each partial sum is
        truncated.
  \item \code{fTimeStamp} = maximum bucket; \code{fTimeStampCorr} = \code{TBCorr};
        \code{fTimeStampCorrInter} = the parabolic sub-bucket offset
        $\frac{1}{2}(a_{-1}-a_{+1})/(a_{-1}+a_{+1}-2a_0)$, stored only, never used for $z$.
  \item \code{SetPeakingShiftTBs(n)} (default 0) subtracts $n$ buckets before the $z$ conversion to
        account for the lag between the electron arrival and the response maximum.
\end{itemize}
In \code{getTBCorr} the loop runs over $i = 0\ldots10$ and tests
\code{adc[max-i+10]} ($>0$ and $<4000$) but sums \code{adc[max-i+5]}: the sample tested and the
sample summed differ by five buckets. When none of the tested samples passes, \code{TBCorr} is 0,
which with \code{SetTimeCorrection(true)} puts the hit at time bucket 0.

Because $z$ is quantised to whole time buckets (2.08~mm in the a1975 parameters), several hits of
a track routinely share exactly the same $z$. Algorithms that assume distinct $z$ values (spline or
ordering by $z$) must handle such ties.

\subsection{\code{AtPSAMultiFit}}

\code{AtPSAMultiFit} fits a sum of GET response pulses, $A\,e^{-3u}\sin(u)\,u^3$ with $u =
(t-t_0)/\tau$, to each trace, one component per seed peak, so that two pulses on the same pad
(two tracks crossing it at different drift times) give two hits.

\begin{center}
\begin{tabular}{p{4.8cm}p{1.6cm}p{8.2cm}}
\toprule
Setter & Default & Meaning \\
\midrule
\code{SetPeakingTime(us)} & 0.720 & Response time constant $\tau$ [\textmu s]. \\
\code{SetMaxPeaks} & 4 & Maximum pulses per pad. \\
\code{SetMinSeparation(tb)} & 8 & Minimum separation of seeds [TB]. \\
\code{SetSeedSigma} & 4.0 & Prominence (in noise RMS) required of additional seeds. \\
\code{SetProminence(adc)} & 0 & If $>0$, absolute prominence required of \emph{all} seeds. \\
\code{SetProminenceOnPrimary} & false & Require prominence of the first seed too. \\
\code{SetMinSignificance} & 3.0 & Keep a pulse only if $A/\sigma_A \ge$ this. \\
\code{SetFloatPeakingTime} & false & Fit a common $\tau$. \\
\code{SetFitChi2Cut} & 0 & If $>0$, reject pulses whose local reduced $\chi^2$ exceeds it (turns on
floating $\tau$). \\
\code{SetFitRelErr} & 0.1 & Relative model error in that $\chi^2$. \\
\code{SetReportT0} & false & Report the impulse arrival $t_0$ instead of the response maximum. \\
\bottomrule
\end{tabular}
\end{center}

The noise RMS is $1.4826\times$ the median absolute deviation of the trace (floor 1 ADC). Seeds are
local maxima of a 3-point smoothed trace in buckets 21--499 above the threshold. The fit range is
from 5 buckets before the first seed to $5\tau$ after the last (at most bucket 510). A pulse is
kept if its amplitude is at least the threshold and its peak time is within 20--500.

Output per pulse: $z$ = \code{CalibrateZ} at the fitted peak time $t_0 + u_\mathrm{peak}\tau$ (a
continuous value, not quantised), charge = fitted amplitude at the peak, trace integral = $A\times
\int\!\mathrm{response}\times\tau$ (with the fixed $\tau$ even when $\tau$ floats),
\code{fChargeVariance} = local reduced $\chi^2$ (0 unless $\tau$ floats), and
\code{fTimeStampCorrInter} = fitted $\tau$. The time stamp is written to the integer
\code{fTimeStamp} and to \code{fTimeStampCorr}.

\paragraph{Saturation.} Nothing in \code{AtPSAMultiFit} identifies or masks clipped samples. The
amplitude parameter is limited to $5A_0+10$, where $A_0$ is derived from the seed height, so the
fit of a clipped pulse is driven by its unclipped edges and can return an amplitude above the ADC
ceiling. The macro \url{macro/Unpack_HDF5/a1975/D2_UKF/saturation_dp.C} histograms hit amplitudes
up to 6000 ADC to look at this.

\subsection{Other PSA classes}

\paragraph{\code{AtPSATBAvg}.} Splits the trace into consecutive windows of \code{SetNumTBToAvg}
buckets (default 5) and makes one hit per window whose mean ADC is at or above the threshold, at
$z$ = \code{CalculateZGeo} of the window centre. Charge = window mean; trace integral = sum of the
whole trace (stored on every hit). \code{SetMaxThreshold(adc)} (default: largest double) skips the
whole pad when its maximum exceeds the value; the code comment calls these pads saturated.
\code{UseArrayAugment(name)} runs on a pad augment (for example a deconvolved charge) instead of
the ADC.

\paragraph{\code{AtPSAHitPerTB}.} One hit per time bucket whose ADC exceeds the threshold, in the
range set by \code{SetTBLimits(\{first, last\})} (default 0--512). Charge = ADC of the bucket. With
\code{SetReplaceTraceIntegral(true)} (default) every hit carries the sum of the above-threshold
samples as trace integral.

\paragraph{\code{AtPSAFull}.} For traces whose longest above-threshold interval is shorter than 50
buckets it makes a single hit at the maximum (buckets 20--499) with charge \emph{zero}: the charge
variable is never assigned in that branch, and the hit is created even when no sample is above
threshold. Longer intervals are split into 50-bucket windows, one hit per window at the window
start, with charge = (sum over the window)/50.

\paragraph{\code{AtPSASpectrum}.} Peak finding with \code{TSpectrum::SearchHighRes}, optional
background suppression and interpolation, optional centre-of-gravity time
(\code{SetTimeCorrection}). As written, the working copy of the trace is filled with zeros
(\code{floatADC.fill(0)}) immediately after being copied from the ADC and before the peak search, so
the search runs on an empty array. Verify that the class produces hits before relying on it.

\paragraph{\code{AtPSASimple2}.} Deprecated (the class attribute recommends \code{AtPSASpectrum} or
\code{AtPSAMax}). It overrides \code{Analyze} directly. One of \code{SetMaxFinder()} or
\code{SetPeakFinder()} must be called: with neither, the number of peaks is zero and no hit is
produced. It has the same tested-versus-summed index offset in the centre-of-gravity time as
\code{AtPSAMax}, an optional baseline correction (\code{SetBaseCorrection}) and gain and jitter
calibration through \code{AtCalibration}.

\paragraph{\code{AtPSADeconv}, \code{AtPSADeconvFit}, \code{AtPSAIterDeconv}.}
\code{AtPSADeconv} divides the Fourier transform of each trace by that of the electronics
response, applies a Butterworth low-pass filter (\code{SetFilterOrder}, even orders only;
\code{SetCutoffFreq}; with the default order 0 no filter is applied), transforms back, subtracts
the mean of the first 20 samples and stores the result as the pad augment \code{"Qreco"}. A response
must be provided with \code{SetResponse}, either as an \code{AtRawEvent} of response pads or as a
function \code{(padNum, time[us])}. \code{AtPSADeconv} itself applies no threshold: every pad gives
one hit at the charge-weighted mean time bucket of the whole deconvolved trace, with $z$ variance
from the time spread. \code{SetUseSimCharge(true)} uses a simulated \code{"Q"} augment when present.
\code{AtPSADeconvFit} instead fits a Gaussian to the deconvolved charge around its maximum
($\pm(3\sigma+2)$ buckets, $\sigma$ from the longitudinal diffusion coefficient \code{CoefL}),
applies the threshold and the 20--500 window, and returns charge $= A\sigma\sqrt{2\pi}$.
\code{AtPSAIterDeconv} adds \code{SetIterations(n)} (default 0) corrections, each deconvolving the
residual between the trace and the reconvolved charge.

\paragraph{\code{AtPSAComposite}.} Constructed as \code{AtPSAComposite(beamPSA, PSA, beamRadius =
40)}: pads whose centre lies within \code{beamRadius} [mm] of the axis are analysed with
\code{beamPSA}, the others with \code{PSA}. Thresholds and options must be set on the two
sub-objects. (The example in \url{docs/subsystems/psa.md} uses an \code{AddPSA} method that does
not exist.)

\subsection{Summary}

\begin{center}
\begin{tabular}{p{3.0cm}p{1.6cm}p{3.6cm}p{3.0cm}p{2.8cm}}
\toprule
Class & Hits per pad & $z$ from & Hit charge & Uses \code{CalibrateZ} \\
\midrule
\code{AtPSAMax} & 1 & integer max TB (or CoG) & peak ADC & yes \\
\code{AtPSAMultiFit} & 1--4 & fitted peak, continuous & fitted amplitude & yes \\
\code{AtPSATBAvg} & many & window centre & window mean & no \\
\code{AtPSAHitPerTB} & many & each TB & ADC & no \\
\code{AtPSAFull} & 1 or many & max TB / window start & 0 / window mean & no \\
\code{AtPSASpectrum} & many & peak TB (or CoG) & peak ADC & no \\
\code{AtPSADeconv} & 1 & CoG of deconvolved charge & integral & no \\
\code{AtPSADeconvFit} & 1 & Gaussian mean & Gaussian area & no \\
\bottomrule
\end{tabular}
\end{center}

%=====================================================================================================
\section{Hit-level cleaning: \code{AtDataCleaningTask}}

\code{AtDataCleaningTask} (\url{AtReconstruction/AtDataCleaningTask.h}) takes a
\code{std::unique\_ptr} to an \code{AtTools::DataCleaning::AtDataCleaner}, reads an \code{AtEvent}
(default \code{"AtEventH"}), and writes a copy whose hits are those returned by
\code{CleanData()} (default output \code{"AtEventCleaned"}, persistent by default). Available
cleaners are \code{AtkNN(k, distance)}, which rejects hits whose mean distance to their $k$ nearest
neighbours exceeds the distance, and \code{AtDirDeDxCleaner}. The simulation macro
\url{macro/Simulation/ATTPC/16C_pd/run_reco_C16pd.C} runs it after the PSA:
\begin{verbatim}
  auto cleaner = std::make_unique<AtTools::DataCleaning::AtDirDeDxCleaner>();
  auto *cleanTask = new AtDataCleaningTask(std::move(cleaner));
  cleanTask->SetInputBranch("AtEventH");
  cleanTask->SetOutputBranch("AtEventClean");
\end{verbatim}

%=====================================================================================================
\section{Worked example: the a1975 data chain}
\label{sec:unp-example}

The macro \url{macro/Unpack_HDF5/a1975/UKF/pipeline/unpackReco_a1975_UKF.C} runs unpacking, PSA,
space-charge correction and pattern recognition on one a1975 run. The parts relevant here are
\begin{verbatim}
  auto fAtMapPtr = std::make_shared<AtTpcMap>();
  fAtMapPtr->ParseXMLMap(mapDir.Data());          // scripts/ANL2023.xml
  fAtMapPtr->GeneratePadPlane();

  auto unpacker = std::make_unique<AtHDFUnpacker>(fAtMapPtr);
  unpacker->SetInputFileName(inputFile.Data()); // /mnt/f/a1975/h5/run_XXXX.h5
  unpacker->SetNumberTimestamps(2);
  unpacker->SetBaseLineSubtraction(true);
  auto unpackTask = new AtUnpackTask(std::move(unpacker));
  unpackTask->SetPersistence(persistRaw);

  auto psa = new AtPSAMax();
  psa->SetThreshold(20);
  AtPSAtask *psaTask = new AtPSAtask(psa); // clones psa: set options first
  psaTask->SetPersistence(persistRaw);
  psaTask->SetOutputBranch("AtEventH");
\end{verbatim}
The parameter file is \url{parameters/ATTPC.a1954.par}. The IC traces come from a separate run of
\url{macro/Unpack_HDF5/a1975/UKF/pipeline/unpackFRIB_a1975_UKF.C} (\code{AtFRIBHDFUnpacker}) and
are paired with this output by entry number, as described in Section~\ref{sec:unp-align}. The older
macro \url{macro/Unpack_HDF5/a1975/unpack_a1975.C} constructs an \code{AtFilterSubtraction}
(threshold 50) and its task but does not add the task to the run, so its PSA reads the unfiltered
\code{AtRawEvent}.

For the deuterium-target runs, \url{macro/Unpack_HDF5/a1975/D2_UKF/unpackReco_multifit.C} uses
\code{AtPSAMultiFit} with the same unpacker settings.

%=====================================================================================================
\section{Pitfalls}
\label{sec:unp-pitfalls}

\begin{enumerate}
  \item \textbf{Scalers are incremental.} Sum every scaler event of every run in the run
        \emph{list} of the analysis (\code{AtFribScalers::ReadRun} and \code{Add}). Do not take the
        last value of a run, do not difference the values at the ends of a range, and do not sum a
        range of run numbers when the analysis used a subset.
  \item \textbf{GET and FRIBDAQ are paired by entry.} Check that both streams start at the same DAQ
        event and have the same number of entries. A truncated auxiliary stream gives no warning,
        and \code{AtDataReductionTask} removes entries from one file only.
  \item \textbf{Baseline subtraction is off by default} in the HDF5 unpackers. Without it pads are
        not flagged pedestal-subtracted and the threshold of the PSA acts on raw ADC values.
  \item \textbf{\code{AtPSAMax} quantises $z$} to whole time buckets; the sub-bucket
        interpolation is only stored in \code{fTimeStampCorrInter}. Hits of a track can share an
        exact $z$, whereas \code{AtPSAMultiFit} gives continuous $z$. A subsequent space-charge
        correction shifts $z$ by a position-dependent amount and removes the exact ties, a
        point the comment in \code{unpackReco\_a1975\_UKF.C} discusses for spline-based
        estimators.
  \item \textbf{Centre-of-gravity index offset.} In \code{AtPSAMax::getTBCorr} (and in
        \code{AtPSASimple2} and \code{AtPSASpectrum}) the sample tested,
        \code{adc[max-i+10]}, is not the sample summed, \code{adc[max-i+5]}.
  \item \textbf{Saturated pulses.} \code{AtPSAMultiFit} does not mask clipped samples and can
        return amplitudes above the ADC ceiling; \code{AtPSATBAvg} with \code{SetMaxThreshold}
        drops the whole pad; \code{AtPSAMax} excludes samples at or above 4000 only from the
        centre-of-gravity time.
  \item \textbf{Same PSA for data and simulation.} Hit charge, $z$ reference (response maximum
        or impulse) and $z$ quantisation all depend on the class and its options. A PSA change on
        one side shifts the $z$ scale and the hit-charge distributions relative to the other.
        As an example, \url{macro/Simulation/ATTPC/16C_pd/run_reco_C16pd.C} labels its
        \code{AtPSAMultiFit} block ``same as the data'', while the a1975 H$_2$ data production
        macro \code{unpackReco\_a1975\_UKF.C} uses \code{AtPSAMax}; confirm which data
        production a simulation is compared with.
  \item \textbf{Raw-pointer constructors clone.} \code{AtPSAtask(AtPSA*)} copies the PSA at
        construction; setters called afterwards have no effect.
  \item \textbf{Inhibit maps.} The HDF5 unpackers ignore \code{AtMap} inhibit flags, while the
        digitization honours them; inhibited pads therefore appear in data but not in
        simulation unless they are removed elsewhere.
  \item \textbf{\code{SetSpyralZ}} uses the opposite $z$ origin from \code{CalculateZGeo};
        combined with the fitters it mirrors $z_\mathrm{lab}$.
  \item \textbf{Event charge} from \code{AtPSA::Analyze} sums the trace integral once per hit,
        so it is inflated for the classes that store the full trace integral on several hits
        per pad (\code{AtPSATBAvg}, \code{AtPSAHitPerTB} with the default option).
\end{enumerate}
